\documentclass[12pt]{article}
\input{bayesuvius.sty}
\usepackage{authblk}
\title{Drought or Flood Warning Decision\\ made via Damage functions and Bayesian Network}
\author[1]{Nishadh Kalladath}
\author[1]{Hillary Koros}
\author[1,2]{Robert R. Tucci}
\affil[1]{ICPAC (IGAD Climate Prediction and Applications Centre)}
\affil[2]{www.ar-tiste.xyz}
\begin{document}
\maketitle
\abstract{xxx}
\section{Introduction}

{\bf Drought/Flood Anticipatory Action} (DFAA)

bnet (Bayesian network) (see Ref.\cite{Bayesuvius})


\begin{figure}[h!]
\centering
\includegraphics[width=6in]
{tim-palmer-cartoon.png}
\caption{Illustration
from Tim Palmer's book Ref.\cite{palmer-book}}
\label{fig-tim-palmer-cartoon}
\end{figure}

Bring the following DFAA models together via
Damage Functions

\begin{enumerate}[(a)]
\item Ref.\cite{palmer-book} book by Tim Palmer

\item Ref.\cite{climada}\label{item-cli} Climada Software

\item Ref.\cite{huizinga} by Huizinga et al.
Global flood depth damage functions

\item Ref.\cite{liu} by Liu et al.
A method of deriving damage fractions functions.
\end{enumerate}


\section{Notation
and Conventions}
We will use $\indi(S)$ 
to denote the indicator
function
(a.k.a. the truth function).
$\indi(S)$ equals 1 if statement $S$ is true,
and 0 if $S$ is false.
For example, $\indi(x\geq 0)$
equals 1 if $x\geq 0$ and 0 otherwise.

In this article, random variables will
be indicated by underlining.
For example, $P(\rvx>x)$
will denote
the probability 
that the random 
variable $\rvx$
is greater than $x$.


\section{Model}

$\alp=$ index that ranges over different hazards


$f_\alp(x_\alp)=$ {\bf Damage Fraction} (term used in Ref.\cite{huizinga}), {\bf Impact Function} (term used in Ref.\cite{climada}) for hazard $\alp$.
$f_\alp(x_\alp)$ is a function of the
parameter $x_\alp$, and increases monotonically
from $0$ when $x_\alp=0$ to $1$ when 
$x_\alp$ is maximum, $x_\alp=x_{\alp, max}$

$D_{\alp, max}=$ {\bf Maximum Damage} (term used in Ref.\cite{huizinga}), {\bf Exposure} (term used in Ref.\cite{climada}) for hazard $\alp$. In dollars.  

Henceforth,
anything called a Damage will
be in units of dollars (or some  other currency),
except damage fractions, which will be unitless,
and range from zero to one.

$D_\alp(x_\alp)=$ {\bf Damage} (term used in Ref.\cite{huizinga}), {\bf Impact} (term used in Ref.\cite{climada})

\beq
D_\alp(x_\alp) = D_{\alp, max} f_\alp(x_\alp)
\eeq

$\ol{D}_\alp=$ {\bf Average Damage}

\beq
\ol{D}_\alp= \frac{1}{{x_{\alp, max}}}\int_0^{x_{\alp, max}} dx\; f_\alp(x_\alp)
\eeq

$D_{max}=$ {\bf maximum total damage}  that 
is  tolerable. 
If total Damage is higher than this, issue drought or flood warning.

We will use $x_\alp={\tt None}$ to indicate

\beq
D_\alp(x_\alp={\tt None}) = \ol{D}_\alp
\eeq


In general, 

$$x_\alp\in val(x_\alp)=\{ {\tt None}, \text{$x_\alp$ for historical events}, \text{bin values of $x_\alp$}\}$$

$\eps_\alp=$ {\bf error} in $D_\alp$,
normally distributed with mean zero
and standard deviation $\s_\alp$.

\beq
D'_\alp(x_\alp) = D_\alp(x_\alp) +\eps_\alp
\eeq

$f_D=$ {\bf total Damage fraction}

\beq
f_D=
\left\{
\begin{array}{ll}
1 & \text{if } \sum_\alp 
D'_\alp(x_\alp)> D_{max}
\\
\frac{1}{D_{max}}\sum_\alp 
D'_\alp(x_\alp)&
\text{otherwise}
\end{array}
\right.
\eeq

Let

$\kappa=(\sigma, \delta)=$
the drought/flood kind, 
where  $\sigma\in\{\text{MAM, JJAS}\}$\footnote{Letters here stand for the first letter of the months of the year.} is the season, and $\delta\in\{\text{Moderate, Extreme, Severe}\}$
is the drought/flood degree.

$S_\kappa=$ {\bf Severity } of drought is defined as minus SPI\footnote{{\bf Standardized Precipitation Index} (SPI) ranges from $-4$ (least rain) to $4$ (most rain).} (higher rain,
higher SPI, less 
severe drought). 
For flood, we would
use $+SPI$ instead.
Severity is called {\bf Intensity}
in Ref.\cite{climada}

$ P(S_\kappa)=$ likelihood, 
probability  that a drought of severity $S_\kappa$  will occur in the near future (weeks to months), based on forecasts.

$\mu=$ ensemble member index

$M=$ number of ensemble members, 51 for $ECMWF$ (European Centre for Medium-Range Weather Forecasts)



\begin{table}[!h]
\begin{tabular}{|l|l|l|l|}
\hline
$\sigma\backslash\delta$ & \cellcolor[HTML]{FFFFC7}Moderate & \cellcolor[HTML]{FFFFC7}Extreme & \cellcolor[HTML]{FFFFC7}Severe \\ \hline
\cellcolor[HTML]{FFFFC7}MAM & $-$0.55 & $-$0.98 & $-$0.99 \\ \hline
\cellcolor[HTML]{FFFFC7}JJAS & $-$0.41 & $-$0.98 & $-$0.99 \\ \hline
\end{tabular}
\caption{SPI threshold  for Karamoja, Uganda.
Cells give $-S^*_\kappa$,
where $\kappa=(\sigma, \delta)$}
\label{tab-spi-karamoja}
\end{table}

$S_\kappa^*  =$ severity threshold 
   for drought (See Table
   \ref{tab-spi-karamoja}
   for examples
   of $S^*_\kappa$.)
   

\beq
P(\rvS_\kappa\geq S_\kappa^*  ) = \frac{1}{M} \sum_{\mu=1}^{M} \indi(S^{\kappa,\mu} \geq S_\kappa^*) 
\eeq

$W_\kappa=$ {\bf drought warning}, either 0 or 1

$f_{D, \kappa}=$ the total damage fraction $f_D$ defined
previously, for drought/flood kind $\kappa$. We 
will use this as a threshold for triggering  a  drought/flood warning.



\beq
W_\kappa=
\indi\left[P(\rvS_\kappa\geq S_\kappa^*  ) > 1-f_{D,\kappa}\right]
\eeq

The above model can be 
expressed via the Bayesian Networks
in Figs.\ref{fig-W-bnet}
and \ref{fig-fd-bnet}

\begin{figure}[h!]
$$
\xymatrix{
\rvP(\rvS_\kappa> S^*_\kappa)\ar[d] 
& \rvf_{D,\kappa}\ar[dl]
\\
\rv{W}_\kappa
}
$$
\caption{bnet for calculating warning $\rvW_\kappa$}
\label{fig-W-bnet}
\end{figure}

\begin{figure}[h!]
$$
\xymatrix{
\rvx_{\alp_1} \ar[d]
& \rveps_{\alp_1}\ar[ld]
&\rvx_{\alp_2} \ar[d]
& \rveps_{\alp_2}\ar[ld]
\\
\rvD'_{\alp_1}(x_{\alp_1})\ar[d]
&&\rvD'_{\alp_2}(x_{\alp_2})\ar[dll]
\\
\rvf_D&\ar[l] \rvD_{max}
}
$$
\caption{bnet for calculating total damage
fraction $f_D$ for two hazards $\alp_1$
and $\alp_2$.}
\label{fig-fd-bnet}
\end{figure}



\bibliographystyle{plain}
\bibliography{aa_references}
\end{document}